\section{Methodology}\label{methodology}

\subsection{Models and the pipeline grid}\label{models-and-the-pipeline-grid}

We study the OLMo-3 7B family \citep{olmo3_2025}, which publishes every stage of
its alignment pipeline. Our grid has 25 model states: the base model
(\texttt{Olmo-3-1025-7B}) and 13 stage-3 pre-training anneal checkpoints (steps
1000--11921), the SFT and DPO snapshots, the final RLVR-tuned Instruct model,
and its 8 intermediate RLVR checkpoints. To these we add one refusal-ablated
Instruct model (Section 3.6). OLMo-3 7B is a 32-layer, 4096-dimensional decoder
with hybrid attention: 24 sliding-window layers and 8 full-attention layers
(layers 3, 7, \dots, 31). For the short probing texts used here every token
attends within the window, so the attention pattern is inert for probing; we
nonetheless flag the full-attention layers in all layer-wise plots and find no
periodicity.

We probe foundation directions on the base model and transfer them across the
grid. All probing uses raw-text inputs: a Sprint-1 comparison found that
base-trained directions transfer to the instruct model with higher and more
uniform agreement under raw text than under the chat template (mean
direction cosine \(0.94\) vs.~\(0.91\)), and raw text is the only format defined
consistently across the base and pre-training checkpoints, which have no chat
template. Behavioral and coupling measurements, which require the model to
respond, use the chat template.

\subsection{Foundation probe directions and transfer}\label{foundation-probe-directions-and-transfer}

Following the geometry method of this series \citep{reblitzrichardson2026geometry}
and standard linear probing \citep{alain2017probes}, we train one linear probe
per Moral Foundations Theory \citep{haidt2012righteous, graham2013mft}
foundation on mean-pooled residual activations of matched moral/neutral sentence
pairs, at each layer, with a fixed seed so the learned weight vector (the
foundation ``direction'') is a deterministic function of the activations. We also
record the mean-difference direction. We report two transfer quantities for a
state: the ROC-AUC of the base direction applied to that state's activations
(does the base direction still separate moral from neutral?), and the cosine
between the base direction and a direction freshly fitted on that state (how
much has the direction rotated?). Bootstrap resampling on the base model
(200 resamples) gives a stable layer band of 15--31 across all foundations
(Appendix B); geometry is reported over that band.

\subsection{Framework geometry}\label{framework-geometry}

From the six foundation directions at a layer we compute the \(6\times6\) cosine
matrix, the mean pairwise cosine (collapse if near 1, isolation if near 0), the
effective dimensionality (principal components for 90\% variance), and a Ward
clustering of \(1-\cos\). Tracking these across the grid measures whether
post-training preserves or restructures the moral representation.

\subsection{Coupling: comprehension vs.~compliance}\label{coupling-comprehension-vs.-compliance}

Comprehension and compliance can each be measured on their own; the question
this paper turns on is how tightly the first predicts the second. We measure
that link per scenario, on the 48 moral scenarios of the behavioral benchmark
(six foundations \(\times\) four difficulty levels: obvious, moderate, subtle,
adversarial), each annotated with a target foundation and an expected
approve/disapprove judgment, presented through the chat template.

\textbf{Comprehension signal (internal).} For a scenario we take the residual-stream
activation at the stable layer (layer 16) at the final prompt token, the
position from which the model begins its answer, and project it onto the six
base-trained foundation directions. The six projections sit on different scales,
so we \(z\)-score each foundation's projection across the 48 scenarios and define
the dominant foundation as the \(\arg\max\) of the \(z\)-scored projections. The
scenario is \emph{comprehended} if its dominant foundation matches the annotated
target. This is a stricter test than probe accuracy: probe accuracy asks whether
a foundation is decodable in aggregate, whereas this asks whether the model
places \emph{this} scenario on the right foundation at the decision point.
Chance is \(1/6 \approx 0.17\). We use the base-trained directions at every state,
so comprehension is read on the same six axes throughout.

\textbf{Compliance signal (behavioral).} We greedily generate a completion and parse
its moral judgment; the scenario \emph{complies} if the judgment matches the
expected one. We use this approve/disapprove judgment, not the comply/refuse
classifier, because the former is parsed reliably from a direct ``is this
acceptable?'\,' prompt while the latter is noisier on borderline requests
(Appendix B).

\textbf{Coupling statistics.} Each scenario yields two bits, comprehended and
complied, and a \(2\times2\) contingency table over the 48 scenarios. We summarize
it three ways: raw agreement (the fraction of scenarios on which the two bits
coincide); the \(\phi\) coefficient (the Pearson correlation between the two binary
variables, equivalently the Matthews correlation of the table), for which
\(\phi = 0\) is statistical independence; and the conditional compliance rates
\(P(\text{comply}\mid\text{comprehend})\) and
\(P(\text{comply}\mid\neg\text{comprehend})\), whose difference is the
predictive-gap form of the same relationship. Reading all three keeps the
agreement number from being mistaken for coupling when it is driven by the base
rates of the two bits.

\textbf{What the measure is and is not.} Coupling here is correlational: it asks
whether the foundation the model internally activates co-varies, across
scenarios, with whether it behaves correctly. It does not show that the moral
representation \emph{causes} the behavior; a causal test would steer the
foundation direction during generation and read the change in judgment, which we
leave to future work. With 48 scenarios and keyword judgment parsing the estimate
is noisy, and difficulty level is a partial confound we do not control, so we
report coupling qualitatively, as the result that the two are barely linked,
rather than as a precise coefficient. Its value is comparative: applying the
same procedure at SFT, DPO, Instruct, and the ablated model makes the trend in
coupling and its small magnitude interpretable even when any single coefficient
is imprecise.

\subsection{Persona direction}\label{persona-direction}

We train a linear persona probe on persona-voice vs.~neutral-voice minimal
pairs (a linear analog of the toxic-persona feature of
\citealp{wang2025persona}) and, per layer, measure its accuracy and the cosine
of the persona direction to each foundation direction (the persona--morality
angle). Tracking both across the grid asks whether alignment couples persona to
morality.

\subsection{Refusal ablation and refusal--morality geometry}\label{refusal-ablation-and-refusalmorality-geometry}

We recover a refusal direction with the Heretic protocol
\citep{pew2025heretic, arditi2024refusal}: its exact prompt set
(\path|mlabonne/harmful_behaviors| and \path|mlabonne/harmless_alpaca|, 400 train prompts
each), the chat template with a generation prompt, and a per-layer
difference-of-means of the final-token residual between harmful and harmless
prompts. We ablate with Arditi et al.'s uniform single-direction method rather
than Heretic's per-layer optimization, a controlled rather than maximal
ablation: the unit refusal direction \(\hat d\) from the stable layer is
orthogonalized out of every layer's attention out-projection and MLP
down-projection, \(W \leftarrow W - \hat d \hat d^\top W\). We then re-run the full
battery (Sections 3.2--3.5) on the ablated model. The refusal--morality geometry
is the cosine of \(\hat d\) to each foundation and the fraction of \(\hat d\)'s norm
lying in the six-foundation span (a least-squares projection): a low fraction
means refusal is carried outside the moral subspace.

Behavioral compliance uses two benchmarks: a moral-scenario judgment benchmark
(per-foundation approve/disapprove accuracy) and a borderline-request refusal
benchmark (the fraction of harmful/borderline requests the model declines). The
refusal classifier flags an opening refusal as a refusal regardless of length;
this corrects a failure mode in which a long ``I'm sorry, but I can't \dots''
decline was scored as compliance (Appendix B).
